% Pista ana-24i-q3 · Anàlisi
% Procedència: PAU Catalunya 2024 · Convocatòria de juny (incidències) · Sèrie 5 · Qüestió 3
\documentclass[11pt,a4paper]{article}
\usepackage{pista}
\pistainfo{Catalunya 2024}{Convocatòria de juny (incidències)}{Sèrie 5}

\begin{document}

\section*{\textcolor{blauPista}{Qüestió 3}}

\noindent Considereu les paràboles $y = f_{a}(x)$, amb $f_{a}(x) = ax^{2} + 2x + 5 - a$, on $a$ és un paràmetre real.

\subsection*{\textit{a)} \normalfont Determineu el valor del paràmetre $a$ per al qual la recta tangent a $y = f_{a}(x)$ en el punt d'abscissa $x = 1$ passa pel punt $(2, 13)$. \hspace*{\fill}\textnormal{[1~punt]}}

\pista{El pendent de la recta tangent que busques és $f_{a}'(1)$, i el punt és $\bigl(1, f_{a}(1)\bigr)$. Per tant:
\[
y - f_{a}(1) = f_{a}'(1)\,(x - 1)
\]
En primer lloc, calcula $f_{a}(1)$ i $f_{a}'(1)$, i substitueix-los a la fórmula. Fixa't que a $f_{a}(1)$ el paràmetre $a$ desapareix; en canvi, a $f_{a}'(1)$ sí que hi surt.

\medskip
A continuació, que la recta passi pel punt $(2, 13)$ vol dir que, si substitueixes $x = 2$ i $y = 13$ a l'equació de la recta tangent, la igualtat es compleix. Obtindràs una equació on l'única incògnita és $a$.}

\subsection*{\textit{b)} \normalfont Calculeu els punts de tall de les paràboles $y = f_{1}(x)$ i $y = f_{3}(x)$. \hspace*{\fill}\textnormal{[0,5~punts]}}

\pista{En primer lloc, escriu les dues paràboles. Fixa't que el paràmetre $a$ apareix dues vegades. Per exemple, per a $a = 1$:
\[
f_{1}(x) = 1 \cdot x^{2} + 2x + 5 - 1 = x^{2} + 2x + 4
\]
Fes el mateix amb $a = 3$.

\medskip
A continuació, iguala $f_{1}(x) = f_{3}(x)$ i resol l'equació: les solucions són les abscisses dels punts de tall. Després, substitueix-les a $f_{1}(x)$, que és la més senzilla, per obtenir les ordenades dels punts de tall. Finalment, la teva resposta han de ser dos punts $P$ i $Q$, i cadascun té dues coordenades.}

\subsection*{\textit{c)} \normalfont Calculeu l'àrea de la regió situada entre les dues paràboles $y = f_{1}(x)$ i $y = f_{3}(x)$. \hspace*{\fill}\textnormal{[1~punt]}}

\pista{Les abscisses dels dos punts de tall de l'apartat \textit{b)} et marquen precisament l'interval sobre el qual integres: $[x_{1}, x_{2}]$. Calcula l'àrea amb el valor absolut de la integral de la diferència:
\[
\text{Àrea} = \left| \int_{x_{1}}^{x_{2}} \bigl(f_{1}(x) - f_{3}(x)\bigr)\,dx \right|
\]
Fes la resta $f_{1}(x) - f_{3}(x)$ abans d'integrar. Obtindràs un polinomi de segon grau, que serà senzill d'integrar.}

\end{document}
